\chapter{Help! Help! I'm being repressed}
%Cooperation, coordination, and conformity

%maybe the chapter is about normative judgement

\section{Notes}
In \textit{Civilization and its Discontents}, Freud argued that, achieving social order requires the sacrifice of personal freedom. For this reason, he saw all rule following as repressive. And yet, we generally follow the rules, even when nobody else is watching and they don't affect anyone else. 

\begin{quote}
    Thomas Hobbes's major contribution to political philosophy is to have shown that the use of force by the state need not be justified entirely through reference to punishment or desert. His analysis of the state of nature is intended to show that private individuals, going about their own business, and bearing no particular ill-will toward one another, will nevertheless become locked in collectively self-defeating interaction patterns. Although they can see that it would be mutually beneficial to help one another, to keep their promises, to refrain from violent predation and to respect each others' liberties, there is no basis for trust between them. They are unable to credibly commit themselves to refraining from acting opportunistically -- accepting the benefits provided by others and yet failing to reciprocate. One way to enhance their credibility is to license some third party to impose a punishment upon them in the event of any compliance failure. \\\citep[891]{Heath2023}
\end{quote}

\begin{quote}
    The major reason that human societies and cultures differ from one another so markedly is that social life is governed by a set of shared norms, or rules of conduct, that are not biologically fixed, but rather are reproduced through social learning. The content of these rules varies enormously from place to place and from time to time. The one thing that does not vary, however, is the fact that these rules are enforced. \\\citep[897]{Heath2023}
\end{quote}

\newpage
Rawls's ``suggestions was that justice be thought of not in the image of divine law, as a set of cosmic principles that apply everywhere and at all times, but that it be seen rather as arising endogenously out of various schemes of cooperation, both constraining and enabling their formation.'' \citep[7]{Heath2022}.

An adult learning a language does not wish to be told that anything goes. 

Coordination is a way to overcome the indeterminacy caused by situations in which my best move depends on your move, and your best move depends on my move.

\section{Conformism}
Humans are culturally conformist by nature. In a way, many organisms are conformist by nature: morning glories generally conform to the pattern of blooming in the morning, bacteria undergo mitosis, and raptors don't wake up one day and decide to try vegetarianism. But these are examples of conforming to biologically determined norms. Humans are the rare creature that -- often thoughtlessly -- copies stuff we see others do, even when it's a totally new behaviour. This conformity bias manifests in following fashion trends, adhering to societal expectations, or even adopting popular opinions.

In 2005, Victoria Horner, a doctoral student studying in Andrew Whiten's psychology lab at St Andrews, published a series of experiments in which children and chimps were given the opportunity to receive a reward by performing a number of steps. In these experiments, an opaque puzzle box was presented with a reward inside -- say, a treat. Before allowing their participants to try to get the rewards, researchers would first demonstrate an overly complex approach to opening the box. They might tap the it in a nonsensical sequence, use an unnecessary tool, or poke certain sides in an irrelevant order. Then it was over to the participants

With these opaque boxes, whose internal mechanics are hidden from view, both two- to six-year-old chimpanzees and three- and four-year-old children were expected to imitate the steps they'd seen to open the box, copying all actions observed, regardless of their apparent causal relevance to the task outcome. The expectation was that, in the absence of visible outcomes, learners would not be able to discern which actions were necessary and which were not, leading to a higher likelihood of imitating all actions performed by the model. And, indeed, this is exactly what they did, not perfectly and not entirely consistently -- they're young -- but generally, they went along with what they'd seen. The chimps aped the demonstration with close to 100\% accuracy, the human children with about 75\%.

A different condition used the same procedures, but this time with clear boxes that made the internal mechanism visible. Here, the expectation was that young participants would selectively emulate only the causally relevant actions, understanding the task's causal structure. And, indeed, this is what was observed in chimpanzees, who generally disregarded irrelevant actions and focused on the essential steps to retrieve the reward. The human children, though, tended to continue to replicate the methods as they were demonstrated, including irrelevant actions, with their fidelity similar to that in the opaque condition.

From this, and many subsequent related studies, researchers have found that human children don't just imitate; they ``over imitate''. While other creatures display incredible intelligence and social learning, our ``over-imitation" tendencies stand apart.  Children meticulously recreate entire sequences, even seemingly pointless components, and they persist in this copying beyond their toddler years. At first glance, this behaviour might seem like a cognitive disadvantage, a failure to distinguish the functional from the frivolous. Yet, this quirk may unlock a crucial  aspect of how human culture came to be so elaborate. Imagine  early humans observing tool-making or fire-starting. Without grasping the why behind each action, faithful duplication was their initial pathway to learning those invaluable skills.

But they don't imitate the demonstrator's breathing rate, nor the hand the demonstrator used (depending on their handedness), the stance taken, or the facial expressions during the demonstration. Observers disregard the exact interval timing between actions, overlook the incidental noises produced by the demonstrator, and feel no obligation to mimic the demonstrator's attire. The manner, timing, and occurrence of eye contact initiated by the demonstrator with the audience are not mirrored. And verbal explanations provided by the demonstrator are left unechoed.

``Not every irregularity indicates a rule violation. One has to know the rule if one wishes to determine whether someone is deviating from it'' \citep[17]{Habermas1987}. And so, as Habermas says, the identity of a rule cannot be reduced to mere patterns of behavior. Instead, it relies on ``intersubjective validity''. This entails that individuals, aware of the norms guiding their actions, can both diverge from and critically evaluate such divergences as breaches of established rules.

Learning as Performance: Perhaps mimicking the act with precision was less about instant efficiency and more about signaling belonging. Demonstrating an ability to perfectly recreate those rituals could have held great  significance within social groups. The very act of "over-imitation" might have signaled  a dedication to learn and the potential to uphold communal practices. The child focused solely on getting the fire to  work might have missed out on subtle cultural cues embedded in the performance surrounding the process.

Shared Language as Ritual: This perspective adds a twist to how we typically imagine language emergence. Early forms of communication may have begun less as  tools to describe immediate needs (i.e. "food over there!"), and instead focused on replicable patterns – rhythmic vocalizations or repeated gestures. Shared rituals strengthened bonds long before they carried precisely understood meaning. These rituals, much like the box experiments, likely involved a degree of "over-imitation". Just as kids didn't simplify the box opening for greater efficiency, those first words or proto-phrases weren't boiled down to their most pragmatic form. Perhaps the intent was to echo the sounds of the tribe with precision.

From Imitation to Innovation: It's vital to note that "over-imitation" does not imply blind conformity forever. This initial copying could provide the crucial building block for more rapid cultural change.  Once a skill or ritual is thoroughly ingrained through mimicry, then individuals within the group have the capacity to experiment and  innovate – adding modifications,  discovering streamlined forms.   With a baseline shared "language" (whether sound-based or symbolic) already established via imitation, those tweaks can rapidly propagate through a broader population.

A Bridge to Hurford: All of this speaks to a core idea present within Hurford's work:  humans became remarkably skilled at cultural transmission, with language as both the beneficiary of this strength and a means to amplify it even further. This cycle of shared ritual, innovation, and teaching may have set us down a unique evolutionary path. And the puzzle  box experiments offer tantalizing hints regarding this journey that began with seemingly mindless conformity.




``The expectation of conformity ordinarily gives everyone a good reason why he himself should conform'' \cite[167]{Lewis1983}.

``In any case in which a language \textit{£} clearly is used by a population \textit{P}, then, it seems that there prevails in \textit{P} a convention of truthfulness and trust in \textit{£}, sustained by an interest in communication'' \cite[168]{Lewis1983}.

Ungrammaticality is good because it increases robustness?

``(1) Some conventions are better, from a payoff standpoint. (2) Some conventions are more likely to emerge. And (3), some conventions are more stable than others. Once they have emerged, they are not likely to be abandoned.'' \citep{OConnor2021b}